\documentclass[10pt]{beamer}

\usetheme[numbering=none]{metropolis}
\usepackage{appendixnumberbeamer}

\usepackage[utf8]{inputenc}

\usepackage{booktabs}
\usepackage[scale=2]{ccicons}

\usepackage{verbatim}

\usepackage{pgfplots}
\usepgfplotslibrary{dateplot}

\usepackage{xspace}
\newcommand{\themename}{\textbf{\textsc{metropolis}}\xspace}


\title{Valaris}
\subtitle{Echtzeit Computergrafik in der Spieleentwicklung}
\date{29. Mai 2018}
\author{Lukas Schölch}
\titlegraphic{\hfill\includegraphics[height=2.5cm]{./Bilder/ValarisLogo.png}}

\begin{document}

\maketitle

\begin{frame}[fragile]{Spielkonzept}

  \begin{itemize}
    \item Futuristisches Weltraumrennspiel für Gelegenheitsspieler
        \begin{itemize}
            \item Lineare Strecken mit Checkpoints
        \end{itemize}
    \item Prozedural generierte Planetenbasisstationen
    \item Energiekugeln die Maximalgeschwindigkeit erhöhen
        \begin{itemize}
            \item Können dem Gegner geklaut werden
            \item Werden an Checkpoints anhand der Position des Spielers wieder aufgefüllt
        \end{itemize}
    \item Powerups auf Strecke einsammelbar
    
  \end{itemize}
  
  
\end{frame}

\begin{frame}{Inspiration}

    \hfil \includegraphics[width=.45\linewidth]{./Bilder/PlanetBase.png}
    \hfil \includegraphics[width=.53\linewidth]{./Bilder/MinimumSpaceStation.jpg}
    \vfil \centering \includegraphics[width=.63\linewidth]{./Bilder/StyleExample.png}

\end{frame}

\begin{frame}[standout]{Entwürfe}

    \begin{center}
        \includegraphics [width=.90\linewidth]{./Bilder/VehicleSelectionMockup.png}
    \end{center}

\end{frame}

\begin{frame}[fragile]{Entwürfe}

  \centering\centering
  \includegraphics [width=1.0\linewidth]{./Bilder/HudMockBackground.png}

\end{frame}

\begin{frame}[fragile]{Anforderungen}
    \begin{itemize}  
        \item InGame - Grafik  
        \begin{itemize}
            \item Physikalisch akurate Darstellung von Materialien
            \item Mehrere dynamsische Lichtquellen mit passenden Schatten
            \item Partikel Effekte
            \item Postprocessing
            \item Frei fliegende Kamera zum Testen
        \end{itemize}
        \pause

        \item Prozedurale Generierung
        \begin{itemize}
            \item Zusammenhängende Strecken mit Start, Ziel und Checkpoints
            \item Deterministisch vom Seed abhängig
            \item Unterschiedliche Streckenführungen, aber keine unmöglichen Abschnitte
            \item Definieren von Extremparametern
        \end{itemize}
        \pause

        \item GUI
        \begin{itemize}
            \item Intuitive Menüführung
            \item Alle Funktionen des Spiels ansteuerbar
            \item HUD
        \end{itemize}
    \end{itemize}
\end{frame}

\begin{frame}{Menüführung}

  \begin{center}
    \begin{tikzpicture}
      \node (img1) {\includegraphics[width=.45\linewidth]{./Bilder/Startmenue.png}};
      \pause
      \node (img2) at (img1.south east) [yshift=1cm] {\includegraphics[width=.45\linewidth]{./Bilder/einstellungen.png}};
      \pause
      \node (img3) at (img2.north east) {\includegraphics[width=.45\linewidth]{./Bilder/pausemenue.png}};
    \end{tikzpicture}
  \end{center}
  
\end{frame}


\begin{frame}{Menüführung}

  \begin{center}
    \begin{tikzpicture}
      \node (img1) {\includegraphics[width=.45\linewidth]{./Bilder/streckenauswahl_einzel.png}};
      \pause
      \node (img2) at (img1.south east) [yshift=1cm] {\includegraphics[width=.53\linewidth]{./Bilder/lobby.png}};
      \pause
      \node (img3) at (img2) {\includegraphics[width=.53\linewidth]{./Bilder/spielinterne_einstellungen.png}};
      \pause
      \node (img4) at (img2.north east) [yshift=0.5cm]{\includegraphics[width=.45\linewidth]{./Bilder/endscreen.png}};
    \end{tikzpicture}
  \end{center}

  \end{frame}

\begin{frame}{HUD}

    \begin{center}
        \includegraphics[width=1.0\linewidth]{./Bilder/HudMockUp.png}
    \end{center}

\end{frame}

\begin{frame}

    \centering
        {\scshape\LARGE Valaris \par}
    
    \vspace{1cm}

    \includegraphics[height=2.5cm]{./Bilder/ValarisLogo.png}

    \vspace{1cm}
    
    \tiny{\url{http://www.4players.de/4players.php/screenshot_list/Allgemein/35673/Screenshots/77064/92491909/Minimum.html}}\par
    \url{https://medium.com/@Jacob_Vo/low-poly-models-are-the-future-626034311c46}\par
    \url{https://www.pinterest.de/pin/444308319480944820/?lp=true}

\end{frame}

%\appendix


\end{document}
